\documentclass[12pt,a4paper]{article}
\usepackage[utf8]{inputenc}
\usepackage[T1]{fontenc}
\usepackage[french]{babel}
\usepackage{geometry}
\usepackage{hyperref}
\usepackage{amsmath,amssymb}
\usepackage{booktabs}
\usepackage{listings}
\usepackage{xcolor}
\usepackage{enumitem}
\geometry{a4paper, margin=2.5cm}

\lstset{
  basicstyle=\ttfamily\small,
  breaklines=true,
  frame=single,
  backgroundcolor=\color{gray!8}
}

\title{Rapport Complet sur l'Expressivité du Langage \texttt{PlanningSpec}}
\author{Analyse technique du DSL de planification}
\date{\today}

\begin{document}

\maketitle
\tableofcontents
\newpage

\section{Introduction}
Ce document présente une analyse complète de l'expressivité du langage \texttt{PlanningSpec}, conçu pour modéliser des problèmes de planification à ressources contraintes. L'objectif est double:
\begin{itemize}
  \item expliquer précisément ce que le langage peut représenter;
  \item documenter, pour chaque contrainte et préférence, la traduction MiniZinc et la logique mathématique sous-jacente.
\end{itemize}

La chaîne technique étudiée est:
\[
\texttt{DSL PlanningSpec} \rightarrow \texttt{AST Langium} \rightarrow \texttt{Générateur MiniZinc} \rightarrow \texttt{Solveur}
\]

\section{Périmètre d'expressivité}
\subsection{Objets modélisés}
Le langage couvre les éléments suivants:
\begin{itemize}
  \item \textbf{Temps discret}: jours et créneaux par jour;
  \item \textbf{Activités}: types d'activité, nombre d'instances, durée;
  \item \textbf{Ressources}: types de ressources et instances nommées;
  \item \textbf{Rôles}: rôles applicables à une activité et type de ressource autorisé;
  \item \textbf{Contraintes dures}: obligatoires pour la faisabilité;
  \item \textbf{Préférences souples}: pénalisées dans une fonction objectif.
\end{itemize}

\subsection{Positionnement}
Le DSL est expressif pour des variantes de planification discrète proches de RCPSP/JSSP simplifiés:
\begin{itemize}
  \item précédences d'activités,
  \item fenêtres temporelles,
  \item contraintes de capacité/exclusivité sur ressources,
  \item affectations imposées/interdites,
  \item optimisation de qualité par pénalité.
\end{itemize}

\section{Modèle mathématique global}
\subsection{Ensembles et paramètres}
Soient:
\begin{align*}
&\mathcal{A} &&\text{: ensemble des instances d'activités (ACT\_INST)} \\
&\mathcal{R} &&\text{: ensemble des ressources (RESOURCE)} \\
&\mathcal{K} &&\text{: ensemble des rôles (ROLE)} \\
&H &&\text{: horizon total en slots (TOTAL\_SLOTS)} \\
&d_i &&\text{: durée de l'activité } i \in \mathcal{A}
\end{align*}

Sous-ensembles typés:
\begin{align*}
&\mathcal{R}_\tau \subseteq \mathcal{R} &&\text{: ressources du type } \tau \\
&\mathcal{A}(a) \subseteq \mathcal{A} &&\text{: instances du type d'activité } a
\end{align*}

\subsection{Variables de décision}
\begin{align*}
&s_i \in \{1,\dots,H\} &&\text{: slot de début de l'activité } i \\
&x_{i,r}\in\{0,1\} &&\text{: la ressource } r \text{ est affectée à } i \\
&y_{i,k,r}\in\{0,1\} &&\text{: la ressource } r \text{ tient le rôle } k \text{ dans } i
\end{align*}

\subsection{Contraintes structurelles générées systématiquement}
\paragraph{Non dépassement journalier}
\[
\forall i\in\mathcal{A}:\quad ((s_i-1)\bmod SLOTS\_PER\_DAY)+d_i \le SLOTS\_PER\_DAY
\]

\paragraph{Non-chevauchement par ressource}
\[
\forall r\in\mathcal{R},\ \forall i<j:\ (x_{i,r}\land x_{j,r}) \Rightarrow (s_i+d_i\le s_j \lor s_j+d_j\le s_i)
\]

\paragraph{Compatibilité rôles/ressources}
Pour chaque activité et rôle:
\begin{align*}
&\text{si le rôle n'est pas applicable: } y_{i,k,r}=0 \\
&y_{i,k,r}=1 \Rightarrow x_{i,r}=1 \\
&\sum_{k} y_{i,k,r}\le 1 \quad \text{(une ressource ne porte pas deux rôles sur la même activité)}
\end{align*}

\section{Contraintes du DSL: sémantique et traduction MiniZinc}
\subsection{\texttt{cardinality\_per\_activity}}
\subsubsection*{Sémantique}
Contrôle le nombre d'affectations par instance d'une activité, soit:
\begin{itemize}
  \item sur un \texttt{target} (type de ressource),
  \item sur un \texttt{role},
  \item ou le cas \texttt{target="slot"}.
\end{itemize}

\subsubsection*{Traduction mathématique}
\paragraph{Cible type de ressource \(\tau\)}
\[
\forall i\in\mathcal{A}(a):\quad
m \le \sum_{r\in \mathcal{R}_\tau} x_{i,r} \le M
\]

\paragraph{Cible rôle \(k\)}
\[
\forall i\in\mathcal{A}(a):\quad
m \le \sum_{r\in \mathcal{R}_{type(k)}} y_{i,k,r} \le M
\]

\paragraph{Cible \texttt{slot}}
Le modèle ayant exactement un début \(s_i\) par instance, la traduction vérifie:
\[
m \le 1 \le M
\]

\subsection{\texttt{resource\_exclusivity}}
\subsubsection*{Sémantique}
Limiter l'utilisation d'une ressource d'un type donné pour une activité donnée selon une portée \texttt{slot} ou \texttt{day}.

\subsubsection*{Traduction}
\paragraph{Scope \texttt{slot}}
\[
\forall r\in\mathcal{R}_\tau,\ \forall t\in\{1,\dots,H\}:\quad
\sum_{i\in\mathcal{A}(a)} \mathbf{1}[x_{i,r}=1 \land s_i\le t < s_i+d_i] \le M
\]

\paragraph{Scope \texttt{day}}
\[
\forall r\in\mathcal{R}_\tau,\ \forall q\in\{1,\dots,NUM\_DAYS\}:\quad
\sum_{i\in\mathcal{A}(a)} \mathbf{1}[x_{i,r}=1 \land dayOf(s_i)=q] \le M
\]

\subsection{\texttt{fixed\_assignment}}
\subsubsection*{Sémantique}
Impose qu'une ressource donnée soit affectée à une instance d'activité (et à un rôle quand défini).

\subsubsection*{Traduction}
\[
x_{i,r}=1,\qquad y_{i,k,r}=1\ \text{(si rôles actifs)}
\]

\subsection{\texttt{forbidden\_assignment}}
\subsubsection*{Sémantique}
Interdit une affectation explicite d'une ressource à un rôle pour une instance.

\subsubsection*{Traduction}
\[
y_{i,k,r}=0
\]

\subsection{\texttt{temporal\_precedence}}
\subsubsection*{Sémantique}
Toute instance de l'activité \texttt{beforeActivity} doit se terminer avant toute instance de \texttt{afterActivity}.

\subsubsection*{Traduction}
\[
\forall i\in\mathcal{A}(a_{before}),\ \forall j\in\mathcal{A}(a_{after}):\quad s_i+d_i \le s_j
\]

\subsection{\texttt{time\_window}}
\subsubsection*{Sémantique}
L'instance doit démarrer et finir dans une fenêtre \([L,U]\).

\subsubsection*{Traduction}
\[
s_i \ge L,\qquad s_i+d_i-1 \le U
\]

\subsection{\texttt{mandatory\_roles}}
\subsubsection*{Sémantique}
Chaque rôle applicable d'une activité doit être rempli au moins une fois par instance.

\subsubsection*{Traduction}
\[
\forall i\in\mathcal{A}(a),\ \forall k\text{ applicable}:\quad \sum_{r\in\mathcal{R}} y_{i,k,r}\ge 1
\]

\subsection{\texttt{instance\_precedence}}
\subsubsection*{Sémantique}
Précédence stricte entre deux instances nommées.

\subsubsection*{Traduction}
\[
s_{i_{before}} + d_{i_{before}} \le s_{i_{after}}
\]

\subsection{\texttt{required\_resource}}
\subsubsection*{Sémantique}
Une ressource donnée doit participer à une instance donnée.

\subsubsection*{Traduction}
\[
x_{i,r}=1
\]

\section{Préférences: modélisation souple et fonction objectif}
\subsection{Principe}
Les préférences ne rendent pas le modèle infaisable. Elles ajoutent des pénalités:
\[
penalty = \sum_{\ell} P_\ell
\]
La résolution est:
\begin{itemize}
  \item \(\min penalty\) s'il existe des préférences;
  \item \texttt{satisfy} sinon.
\end{itemize}

\subsection{\texttt{avoid\_participation\_on\_date}}
\subsubsection*{Sémantique}
Éviter qu'une ressource participe à au moins une activité sur un jour donné.

\subsubsection*{Traduction}
Si la date correspond au jour \(q\):
\[
P = w\cdot \mathbf{1}\left[\sum_{i\in\mathcal{A}} \mathbf{1}[x_{i,r}=1 \land dayOf(s_i)=q] > 0\right]
\]

\subsection{\texttt{max\_per\_scope}}
\subsubsection*{Sémantique}
Préférence de plafonnement souple du nombre d'affectations au-delà d'un maximum.

\subsubsection*{Traduction (scope \texttt{slot})}
\[
P = \sum_{r\in\mathcal{R}_\tau}\sum_{t=1}^{H}
w\cdot \max\left(0,\ \sum_{i\in\mathcal{A}(a)}\mathbf{1}[x_{i,r}=1 \land s_i\le t<s_i+d_i]-M\right)
\]

\subsubsection*{Traduction (scope \texttt{day})}
\[
P = \sum_{r\in\mathcal{R}_\tau}\sum_{q=1}^{NUM\_DAYS}
w\cdot \max\left(0,\ \sum_{i\in\mathcal{A}(a)}\mathbf{1}[x_{i,r}=1 \land dayOf(s_i)=q]-M\right)
\]

\subsection{\texttt{preferred\_resource}}
\subsubsection*{Sémantique}
Favoriser une ressource donnée pour un couple (activité, rôle) ou une activité sans rôles.

\subsubsection*{Traduction}
\[
P=
\begin{cases}
w\cdot \mathbf{1}[y_{i,k,r}=0] & \text{si rôles actifs}\\
w\cdot \mathbf{1}[x_{i,r}=0] & \text{sinon}
\end{cases}
\]

\section{Lien explicite entre DSL et fragments MiniZinc}
\subsection{Exemple de solve}
\begin{lstlisting}
solve :: int_search([start_time[i] | i in ACT_INST],
                    first_fail, indomain_min, complete)
       minimize penalty;
\end{lstlisting}

\subsection{Lecture logique}
\begin{itemize}
  \item \texttt{first\_fail}: instancie d'abord les variables les plus contraintes;
  \item \texttt{indomain\_min}: essaye les valeurs basses en premier;
  \item \texttt{complete}: recherche complète.
\end{itemize}

\section{Évaluation de robustesse et d'expressivité}
\subsection{Forces}
\begin{itemize}
  \item séparation nette dur/souple;
  \item couverture riche des contraintes métiers usuelles;
  \item traduction MiniZinc structurée et traçable;
  \item langage éditable facilement (syntaxe JSON-like).
\end{itemize}

\subsection{Limites}
\begin{itemize}
  \item \texttt{temporal\_precedence} est actuellement de type all-to-all;
  \item \texttt{target="slot"} dans \texttt{cardinality\_per\_activity} est une vérification de cohérence plutôt qu'une décision;
  \item pas d'objectif makespan activé par défaut (choix de performance);
  \item validation sémantique encore extensible (cohérence globale, cycles, etc.).
\end{itemize}

\subsection{Conclusion}
Le langage est aujourd'hui suffisamment robuste pour des scénarios de planification discrète contraints orientés production, tout en restant extensible vers des formulations plus avancées (multi-objectif, critères hiérarchiques, variantes RCPSP/JSSP plus strictes).

\section{Annexe A: Récapitulatif synthétique des traductions}
\begin{center}
\begin{tabular}{@{}p{4.2cm}p{9.5cm}@{}}
\toprule
\textbf{Élément DSL} & \textbf{Traduction mathématique principale} \\
\midrule
\texttt{cardinality\_per\_activity} &
\(m \le \sum x_{i,r} \le M\) ou \(m \le \sum y_{i,k,r} \le M\) \\
\texttt{resource\_exclusivity} &
\(\sum \mathbf{1}[\text{actif au slot/jour}] \le M\) \\
\texttt{fixed\_assignment} &
\(x_{i,r}=1\), éventuellement \(y_{i,k,r}=1\) \\
\texttt{forbidden\_assignment} &
\(y_{i,k,r}=0\) \\
\texttt{temporal\_precedence} &
\(s_i+d_i \le s_j\) pour \(i\in A_{before}, j\in A_{after}\) \\
\texttt{time\_window} &
\(L \le s_i\), \(s_i+d_i-1 \le U\) \\
\texttt{mandatory\_roles} &
\(\sum_r y_{i,k,r}\ge 1\) \\
\texttt{instance\_precedence} &
\(s_{i_b}+d_{i_b}\le s_{i_a}\) \\
\texttt{required\_resource} &
\(x_{i,r}=1\) \\
\texttt{avoid\_participation\_on\_date} &
\(w\cdot \mathbf{1}[\exists\ i: x_{i,r}=1 \land dayOf(s_i)=q]\) \\
\texttt{max\_per\_scope} &
\(w\cdot \max(0,\text{charge}-M)\) agrégé par slot/jour \\
\texttt{preferred\_resource} &
\(w\cdot \mathbf{1}[\text{ressource préférée non choisie}]\) \\
\bottomrule
\end{tabular}
\end{center}

\end{document}

